%======================================================================
\chapter{Sampling Profiler}\label{ch:profiler}
%======================================================================

MScheme includes a built-in sampling profiler that can be used to
identify performance bottlenecks in Scheme programs.  The profiler
operates by running a background thread that periodically samples
which procedure is currently executing and tallying the results.
In addition to the flat profile, the profiler records caller/callee
arcs, providing a call graph similar to that produced by {\tt gprof}.

\section{Overview}

Traditional instrumentation-based profilers record the time spent in
every procedure call.  This is accurate but imposes significant
overhead on every call, which can distort the very measurements one is
trying to make.  A~sampling profiler takes a different approach: it
periodically inspects which procedure is currently running and
increments a counter for that procedure.  Given enough samples, the
resulting histogram converges to the true time distribution.  The
overhead is negligible because the interpreter's hot path does only
two pointer writes per procedure call (the current procedure name and
the caller name); all the bookkeeping is confined to the sampler
thread, which wakes up independently on a timer.

\section{API Reference}

\subsection{Flat Profile}

\begin{description}

\item[{\tt (sampling-profiler-start)} / {\tt (sampling-profiler-start {\it interval-ms})}]
  Start the sampling profiler.  The optional argument {\it interval-ms\/}
  specifies the sampling interval in milliseconds; it defaults to~1.
  If the profiler is already running, this procedure has no effect.
  Returns {\tt \#t}.

\item[{\tt (sampling-profiler-stop)}]
  Stop the sampling profiler.  Blocks until the sampler thread has
  exited.  The accumulated samples are retained and can still be
  inspected.  Returns {\tt \#t}.

\item[{\tt (sampling-profiler-reset)}]
  Clear all accumulated samples (both flat profile and call graph) and
  reset the total count to zero.
  Returns {\tt \#t}.  This does not stop a running profiler; to start a
  fresh profiling session one typically calls:
\begin{lstlisting}[language=Scheme]
(sampling-profiler-reset)
(sampling-profiler-start)
\end{lstlisting}

\item[{\tt (sampling-profiler-results)}]
  Return the profiling data as an association list sorted in descending
  order by sample count.  Each element is a pair
  {\tt ({\it name\/}~.~{\it count\/})}, where {\it name\/} is a string
  naming the procedure and {\it count\/} is the number of times the
  sampler observed that procedure as the current one.  If no samples
  have been collected, the empty list is returned.
\begin{lstlisting}[language={}]
(sampling-profiler-results)
=> (("<" . 2027) ("-" . 1763) ("fib" . 1354) ("+" . 452))
\end{lstlisting}

\item[{\tt (sampling-profiler-total)}]
  Return the total number of samples collected as an exact integer.

\item[{\tt (sampling-profiler-report)} / {\tt (sampling-profiler-report {\it n})}]
  Print a formatted profiler report to the current output port.  The
  optional argument~{\it n\/} limits the output to the top~{\it n\/}
  procedures by sample count; by default all procedures are shown.
  Returns {\tt \#t}.
\begin{lstlisting}[language={}]
(sampling-profiler-report)

  samples   %     procedure
  -------  ----   ---------
     2027   36.2   <
     1763   31.5   -
     1354   24.2   fib
      452    8.1   +
  -------
     5596 total samples
\end{lstlisting}

\end{description}

\subsection{Call Graph}

\begin{description}

\item[{\tt (sampling-profiler-call-graph)}]
  Return the call graph as an association list.  Each element has the
  form {\tt ({\it callee\/} ({\it caller\/}~.~{\it count\/}) \ldots)},
  where {\it callee\/} is the procedure that was observed running, and
  each {\tt ({\it caller\/}~.~{\it count\/})} pair indicates how many
  times {\it caller\/} was the most recent call site when {\it callee\/}
  was sampled.  The entries are sorted by total callee sample count,
  descending.  This is the same information that {\tt gprof} uses to
  construct its call graph.

\begin{lstlisting}[language={}]
(sampling-profiler-call-graph)
=> (("<"   ("-" . 2089) ("<" . 3))
    ("-"   ("-" . 3) ("fib" . 1753))
    ("fib" ("+" . 623) ("<" . 643) ("fib" . 9) ("-" . 1))
    ("+"   ("<" . 398) ("+" . 2)))
\end{lstlisting}

  Reading this output for the {\tt fib} example: {\tt fib} was sampled
  1276~times in total.  Of those, 643~samples had {\tt <} as the
  caller (from the test {\tt (< n 2)} calling back into {\tt fib} via
  the {\tt if}), and 623~had {\tt +} (from {\tt (+ (fib \ldots) (fib
  \ldots))}).

\end{description}

\subsection{Call Graph Report}

The {\tt profiling} library (loaded via {\tt (require-modules
"profiling")}) provides a formatted call graph report:

\begin{description}

\item[{\tt (sampling-profiler-call-graph-report)} / {\tt (sampling-profiler-call-graph-report {\it n})}]
  Print a formatted call graph report to the current output port,
  showing each procedure with its callers indented below it.
  The optional argument~{\it n\/} limits the output to the
  top~{\it n\/} procedures.  Returns {\tt ok}.

\begin{lstlisting}[language={}]
(require-modules "profiling")
(sampling-profiler-call-graph-report)

  samples   %     procedure
  -------  ----   ---------
     2148 38.4%   <
                2145  <- -
                   2  <- <
                   1  <- +

     1705 30.5%   -
                1695  <- fib
                  10  <- -

     1313 23.4%   fib
                 671  <- +
                 633  <- <
                   9  <- fib

      423  7.5%   +
                 419  <- <
                   4  <- +

  -------
     5589 total samples
\end{lstlisting}

  Each top-level line shows a callee with its flat sample count and
  percentage.  The indented lines below show how many of those samples
  had each caller as the most recent call site.  For instance, {\tt <}
  was sampled 2148~times, and in 2145 of those the caller was~{\tt -}
  --- reflecting the pattern {\tt (< n 2)} where {\tt n} was just
  produced by a subtraction.

\end{description}


\section{Examples}

\subsection{Profiling a Recursive Function}

The simplest use of the profiler is to bracket a computation with
{\tt start} and {\tt stop} calls:

\begin{lstlisting}[language=Scheme]
(require-modules "profiling")

(define (fib n)
  (if (< n 2) n
      (+ (fib (- n 1)) (fib (- n 2)))))

(sampling-profiler-reset)
(sampling-profiler-start)
(fib 35)
(sampling-profiler-stop)
(sampling-profiler-report)
(sampling-profiler-call-graph-report)
\end{lstlisting}

\noindent
The flat report reveals that the naive Fibonacci spends most of its
time in the arithmetic primitives {\tt <}, {\tt -}, and~{\tt +}, with
the closure {\tt fib} itself accounting for roughly a quarter of the
samples.  The call graph report shows how these procedures relate:
{\tt fib} is called primarily from {\tt +} and {\tt <}, while {\tt -}
is called almost exclusively from {\tt fib}.

\subsection{Comparing Two Implementations}

The profiler can help compare alternative implementations of the same
algorithm.  For instance, one might compare a naive list-based {\tt
  member} against the built-in:

\begin{lstlisting}[language=Scheme]
(define (my-member x lst)
  (cond ((null? lst) #f)
        ((equal? x (car lst)) lst)
        (else (my-member x (cdr lst)))))

(define data (let loop ((i 0) (acc '()))
               (if (= i 10000) acc
                   (loop (+ i 1) (cons i acc)))))

;; Profile the custom version
(sampling-profiler-reset)
(sampling-profiler-start)
(do ((i 0 (+ i 1))) ((= i 100))
  (my-member 5000 data))
(sampling-profiler-stop)
(sampling-profiler-report 5)

;; Profile the built-in
(sampling-profiler-reset)
(sampling-profiler-start)
(do ((i 0 (+ i 1))) ((= i 100))
  (member 5000 data))
(sampling-profiler-stop)
(sampling-profiler-report 5)
\end{lstlisting}

\noindent
The first report will show samples distributed across {\tt
  my-member}, {\tt equal?}, {\tt null?}, {\tt car}, and {\tt cdr}.
The second report will show samples concentrated in the built-in {\tt
  member}, which performs the same work in a single Modula-3 loop
without re-entering the interpreter on each recursive step.

\section{Design Notes}

The sampling profiler is implemented in two parts.  In the eval loop
({\tt Scheme.m3}), every procedure call writes the name of the callee
and the name of the caller (the previously current procedure) to a
pair of global variables.  These writes are guarded by a boolean flag
and cost two pointer-sized stores when profiling is enabled, or one
boolean test when it is not.

A separate Modula-3 thread, forked by {\tt sampling-profiler-start},
wakes up at the configured interval, reads the current procedure and
caller names, and updates two data structures protected by a mutex: a
flat count table (procedure name to sample count) and a call graph
table (callee name to a table of caller names and edge counts).  The
mutex is never acquired by the eval loop, so profiling does not
introduce lock contention in the interpreter.  Both tables key on
text value (not object identity), so samples for the same procedure
are correctly aggregated regardless of which interpreter thread
produced them.

Both closures and primitive procedures are tracked.  Closures receive
their names from {\tt define}; anonymous closures appear as
``anonymous procedure''.

The caller/callee relationship mirrors the approach taken by {\tt
  gprof}: each sample records a single arc from the most recent caller
to the current callee.  For tail calls, the caller is the procedure
that initiated the tail-call chain, which is the correct attribution
--- {\tt f} tail-calling {\tt g} which tail-calls {\tt h} produces
arcs {\tt f}$\to${\tt g} and {\tt g}$\to${\tt h}, exactly reflecting
the actual call sequence.
